\lecture{5}{Thu 24 Sep 2026 09:46}{CW complexes and $\pi _1$}

\begin{proposition}\label{prp:pi1-cw}
	Let $X$ be a space.
	\begin{enumerate}
		\item If $Y$ is obtained by attaching 2-cells to $X$, then the inclusion $X\subseteq Y$ induces a surjection $\pi _1(X,x_0)\to \pi _1(Y,y_0)$ whose kernel is generated by elements $\gamma _\alpha \varphi _\alpha \gamma ^{-1} _\alpha $, for $\gamma $ some path $x_0 \to s$ with $s\in S^2$ a basepoint, and $\varphi _\alpha $ the attaching map.
		\item If $Y$ is obtained by attaching $n$-cells to $X$ for $n>2$, then $X\subseteq Y$ induces an isomorphism $\pi _1(X,x_0) \cong \pi _1(Y,y_0)$.
		\item If $X$ is a path-connected CW complex, then the inclusion of the 2-skeleton induces an isomorphism $\pi _1(X^2,x_0) \subseteq \pi _1(X,x_0)$.
	\end{enumerate}
\end{proposition}


\section{Covering spaces}


\begin{definition}\label{dfn:covering-space-category}
Recall \cref{dfn:covering-space}. Given an object $X$, we define a category of covering spaces of $X$. Given covering spaces $A,B$ of $X$, a map $A\to B$ is defined as a continuous map which commutes with the covering maps. In other words, a map $\varphi $ such that the diagram
\[
	\begin{tikzcd}[row sep = 2.5em, column sep = 2.5em]
	A \ar[" p_A "', dr] \ar[" \varphi  ", rr]  &  & B \ar[" p_B ", dl]  \\
	 & X & 
\end{tikzcd}
\]
commutes. The isomorphisms are precisely when $\varphi $ is a homeomorphism. Automorphisms are also referred to as \emph{deck transformations}. We do not often require a notation for this category, but we may refer to it as $\mathcal{C}\mathrm{over} (X) $.
\end{definition}

\begin{remark}
	With \cref{dfn:covering-space} as Hatcher gave it, the covering map $p$ need not be surjective. Let $x\in X$, and let $x\in U\subseteq X$ be open. Then the preimage could be empty. By taking the index set $J$ to be empty, we have vacuously that for all $\alpha \in J$, the restriction $p|\widetilde{U}_\alpha \cong U$ is a homeomorphism. This is valid, since we require
	\[
		p ^{-1} (U) = \coprod_{\alpha \in J} \widetilde{U} _\alpha ,
	\]
	and the empty disjoint union is the colimit over 0 which is the initial object, hence $\varnothing $ in $\Top $.
\end{remark}

